\documentclass[11pt,a4paper]{article}
\usepackage[T1]{fontenc}
\usepackage{lmodern,array,amsmath,amssymb,booktabs,enumitem,fancyhdr,lastpage,needspace,microtype}
\usepackage[margin=19mm,headheight=14pt,footskip=11mm]{geometry}
\usepackage[hidelinks]{hyperref}
\setlength{\parindent}{0pt}
\setlength{\parskip}{5pt}
\setlength{\emergencystretch}{2em}
\pagestyle{fancy}\fancyhf{}
\fancyhead[L]{\small CS182 --- Lecture 03}
\fancyhead[R]{\small Personal answer copy}
\fancyfoot[L]{\small Independent practice quiz}
\fancyfoot[R]{\small Page \thepage\ of \pageref{LastPage}}
\newcommand{\SetAnswer}[2]{\expandafter\gdef\csname QuizAnswer#1\endcsname{#2}}
\newcommand{\GetAnswer}[1]{\csname QuizAnswer#1\endcsname}
\newlength{\BlankAnswerHeight}\setlength{\BlankAnswerHeight}{16mm}
\newcommand{\AnswerBlank}{\leavevmode\par\vspace{\BlankAnswerHeight}}
% 本文件为独立可填写答题卷；用 pdfLaTeX 编译两次。
% 在 BEGIN TF ANSWERS 和 BEGIN MC ANSWERS 中填写短答案。
% 大题请搜索 BEGIN ANSWER Q25(a) 等注释；每个小问都有独立区域。
% 写完小问后删除该处的 \AnswerBlank，或将 BlankAnswerHeight 改为 0mm。
% 所有打印内容及你填写的回答请使用英文。中文只出现在源文件注释中。
% build_per_lecture.py 不会覆盖已经存在的 Quiz.tex，以保留你的作答。
% BEGIN STUDENT INFO / 个人信息
\newcommand{\StudentName}{}
\newcommand{\StudentID}{}
% END STUDENT INFO
% BEGIN TF ANSWERS / 判断题答题区（第二组花括号填 T 或 F）
\SetAnswer{TF01}{} % Question 1
\SetAnswer{TF02}{} % Question 2
\SetAnswer{TF03}{} % Question 3
\SetAnswer{TF04}{} % Question 4
\SetAnswer{TF05}{} % Question 5
\SetAnswer{TF06}{} % Question 6
\SetAnswer{TF07}{} % Question 7
\SetAnswer{TF08}{} % Question 8
\SetAnswer{TF09}{} % Question 9
\SetAnswer{TF10}{} % Question 10
\SetAnswer{TF11}{} % Question 11
\SetAnswer{TF12}{} % Question 12
% END TF ANSWERS
% BEGIN MC ANSWERS / 不定项选择答题区（填选项字母，如 BD；仅格式示例）
\SetAnswer{MC13}{} % Question 13
\SetAnswer{MC14}{} % Question 14
\SetAnswer{MC15}{} % Question 15
\SetAnswer{MC16}{} % Question 16
\SetAnswer{MC17}{} % Question 17
\SetAnswer{MC18}{} % Question 18
\SetAnswer{MC19}{} % Question 19
\SetAnswer{MC20}{} % Question 20
\SetAnswer{MC21}{} % Question 21
\SetAnswer{MC22}{} % Question 22
\SetAnswer{MC23}{} % Question 23
\SetAnswer{MC24}{} % Question 24
% END MC ANSWERS

\begin{document}
\begin{center}{\Large\bfseries CS182: Introduction to Machine Learning}\par{\Large\bfseries Lecture 03: Probability and Statistics}\par{\large Lecture Quiz and Answer Copy}\end{center}
\textbf{Scope:} Probability\_Statistics.pdf. Independent practice material based on the official 2026 Fall lecture.\par\textbf{Total:} 100 marks. \quad \textbf{Suggested time:} 120 minutes.
\par\textbf{Name:} \underline{\makebox[60mm][l]{\StudentName}}\hfill\textbf{Student ID:} \underline{\makebox[45mm][l]{\StudentID}}
\begin{center}\begin{tabular}{lrrr}\toprule
Section & Questions & Each & Marks\\\midrule
A. Misconceptions: True/False & 12 & 1 & 12\\
B. Multiple response & 12 & 2 & 24\\
C. Proofs and formula derivations & 4 & 10 & 40\\
D. Applications & 2 & 12 & 24\\\midrule
Total & 30 & & 100\\\bottomrule\end{tabular}\end{center}
\textbf{Scoring.} In Section B each item has exactly four options, of which one,
two, or three are correct. The complete correct set earns 2 marks; a nonempty
proper subset earns 1 mark; any incorrect option or no selection earns 0.
T/F answers earn 1 mark only when correct. Show reasoning in Sections C and D;
part marks are available for valid intermediate work. Equivalent correct methods receive credit.
\par\medskip\textbf{Instructions.} Answer in English. Use natural logarithms.
Exact values and justified approximations are accepted. All vectors are columns
unless stated otherwise. Use local assumptions, label conventions, and tie rules.
\par\medskip\textbf{Answer sheet --- Section A}
\begin{center}\renewcommand{\arraystretch}{1.15}\begin{tabular}{|c|c|c|c|c|c|}
\hline
1 & 2 & 3 & 4 & 5 & 6\\\hline
\makebox[20mm]{\rule{0pt}{5mm}\GetAnswer{TF01}} & \makebox[20mm]{\rule{0pt}{5mm}\GetAnswer{TF02}} & \makebox[20mm]{\rule{0pt}{5mm}\GetAnswer{TF03}} & \makebox[20mm]{\rule{0pt}{5mm}\GetAnswer{TF04}} & \makebox[20mm]{\rule{0pt}{5mm}\GetAnswer{TF05}} & \makebox[20mm]{\rule{0pt}{5mm}\GetAnswer{TF06}}\\\hline
\end{tabular}\end{center}
\begin{center}\renewcommand{\arraystretch}{1.15}\begin{tabular}{|c|c|c|c|c|c|}
\hline
7 & 8 & 9 & 10 & 11 & 12\\\hline
\makebox[20mm]{\rule{0pt}{5mm}\GetAnswer{TF07}} & \makebox[20mm]{\rule{0pt}{5mm}\GetAnswer{TF08}} & \makebox[20mm]{\rule{0pt}{5mm}\GetAnswer{TF09}} & \makebox[20mm]{\rule{0pt}{5mm}\GetAnswer{TF10}} & \makebox[20mm]{\rule{0pt}{5mm}\GetAnswer{TF11}} & \makebox[20mm]{\rule{0pt}{5mm}\GetAnswer{TF12}}\\\hline
\end{tabular}\end{center}
\par\medskip\textbf{Answer sheet --- Section B}
\begin{center}\renewcommand{\arraystretch}{1.15}\begin{tabular}{|c|c|c|c|c|c|}
\hline
13 & 14 & 15 & 16 & 17 & 18\\\hline
\makebox[20mm]{\rule{0pt}{5mm}\GetAnswer{MC13}} & \makebox[20mm]{\rule{0pt}{5mm}\GetAnswer{MC14}} & \makebox[20mm]{\rule{0pt}{5mm}\GetAnswer{MC15}} & \makebox[20mm]{\rule{0pt}{5mm}\GetAnswer{MC16}} & \makebox[20mm]{\rule{0pt}{5mm}\GetAnswer{MC17}} & \makebox[20mm]{\rule{0pt}{5mm}\GetAnswer{MC18}}\\\hline
\end{tabular}\end{center}
\begin{center}\renewcommand{\arraystretch}{1.15}\begin{tabular}{|c|c|c|c|c|c|}
\hline
19 & 20 & 21 & 22 & 23 & 24\\\hline
\makebox[20mm]{\rule{0pt}{5mm}\GetAnswer{MC19}} & \makebox[20mm]{\rule{0pt}{5mm}\GetAnswer{MC20}} & \makebox[20mm]{\rule{0pt}{5mm}\GetAnswer{MC21}} & \makebox[20mm]{\rule{0pt}{5mm}\GetAnswer{MC22}} & \makebox[20mm]{\rule{0pt}{5mm}\GetAnswer{MC23}} & \makebox[20mm]{\rule{0pt}{5mm}\GetAnswer{MC24}}\\\hline
\end{tabular}\end{center}
\newpage\section*{Section A. True/False \hfill 12 marks}
\par\noindent\begin{minipage}{\linewidth}
\textbf{1.} If two events have positive probabilities and are mutually exclusive, then they are independent.
\end{minipage}\par\vspace{3mm}
\par\noindent\begin{minipage}{\linewidth}
\textbf{2.} For iid nonnormal observations with finite positive variance, the central limit theorem makes the standardized sample mean exactly standard normal at every finite sample size.
\end{minipage}\par\vspace{3mm}
\par\noindent\begin{minipage}{\linewidth}
\textbf{3.} If \(T\sim t_3\), its variance is one because it is a standardized ratio.
\end{minipage}\par\vspace{3mm}
\par\noindent\begin{minipage}{\linewidth}
\textbf{4.} If two variables each have a normal marginal distribution and have zero covariance, they must be independent.
\end{minipage}\par\vspace{3mm}
\par\noindent\begin{minipage}{\linewidth}
\textbf{5.} The identity \(E(X+Y)=E(X)+E(Y)\) requires independence.
\end{minipage}\par\vspace{3mm}
\par\noindent\begin{minipage}{\linewidth}
\textbf{6.} The weak law of large numbers says that, for a sufficiently large finite sample, the sample mean equals the population mean with probability one.
\end{minipage}\par\vspace{3mm}
\par\noindent\begin{minipage}{\linewidth}
\textbf{7.} For iid observations with finite variance \(\sigma^2\), the variance of their sample mean is \(\sigma^2/n\).
\end{minipage}\par\vspace{3mm}
\par\noindent\begin{minipage}{\linewidth}
\textbf{8.} For a random variable with a probability density, the density value \(p(x)\) may exceed one even though \(P(X=x)=0\).
\end{minipage}\par\vspace{3mm}
\par\noindent\begin{minipage}{\linewidth}
\textbf{9.} For finite-variance variables, \(\operatorname{Var}(X-Y)=\operatorname{Var}(X)+\operatorname{Var}(Y)-2\operatorname{Cov}(X,Y)\).
\end{minipage}\par\vspace{3mm}
\par\noindent\begin{minipage}{\linewidth}
\textbf{10.} For iid \(N(\mu,\sigma^2)\) observations with \(n\ge2\), \((n-1)S^2/\sigma^2\sim\chi^2_{n-1}\), where \(S^2=\sum_i(X_i-\bar X)^2/(n-1)\).
\end{minipage}\par\vspace{3mm}
\par\noindent\begin{minipage}{\linewidth}
\textbf{11.} If \(Z_1,\ldots,Z_k\) are independent standard normal variables, then \(\sum_i Z_i^2\) has a chi-square distribution with \(k\) degrees of freedom.
\end{minipage}\par\vspace{3mm}
\par\noindent\begin{minipage}{\linewidth}
\textbf{12.} For events with \(P(A),P(B)>0\), the identity \(P(A\mid B)=P(B\mid A)\) holds even when the prior probabilities differ.
\end{minipage}\par\vspace{3mm}
\newpage\section*{Section B. Multiple Response \hfill 24 marks}
\textit{Select all correct options. Each question has 1--3 correct options.}\par
\par\noindent\begin{minipage}{\linewidth}
% BEGIN QUESTION Q13 / 题目与选项区
\textbf{13.} Let \(X_1,\ldots,X_n\) be iid with mean \(7\) and finite variance \(9\). Select every justified conclusion.
\begin{enumerate}[label=(\Alph*),leftmargin=8mm,itemsep=2pt,topsep=2pt]
\item For every finite \(n\), \(\sum_iX_i\sim N(7n,9n)\).
\item \(P(|\bar X_n-7|>0.1)\to0\) as \(n\to\infty\).
\item \(\operatorname{Var}(\bar X_n)=9/n^2\).
\item For large \(n\), \(\bar X_n\) is approximately \(N(7,9/n)\).
\end{enumerate}
% END QUESTION Q13

\end{minipage}\par\vspace{3mm}
\par\noindent\begin{minipage}{\linewidth}
% BEGIN QUESTION Q14 / 题目与选项区
\textbf{14.} Select every valid probability identity. Assume all conditional probabilities displayed are defined.
\begin{enumerate}[label=(\Alph*),leftmargin=8mm,itemsep=2pt,topsep=2pt]
\item \(P(A\cap B\cap C)=P(A)P(B\mid A)P(C\mid B)\).
\item \(P(A\mid B)=P(A\cap B)/P(A)\).
\item If \(A_1,A_2\) partition the sample space, then \(P(B)=P(B\mid A_1)+P(B\mid A_2)\).
\item \(P(A\cup B)=P(A)+P(B)-P(A\cap B)\).
\end{enumerate}
% END QUESTION Q14

\end{minipage}\par\vspace{3mm}
\par\noindent\begin{minipage}{\linewidth}
% BEGIN QUESTION Q15 / 题目与选项区
\textbf{15.} A joint density is \(p(x,y)=2\) on \(0<x<y<1\), and zero elsewhere. Select every correct statement.
\begin{enumerate}[label=(\Alph*),leftmargin=8mm,itemsep=2pt,topsep=2pt]
\item \(p_X(x)=2(1-x)\) for \(0<x<1\).
\item \(X\) and \(Y\) are independent because their joint density is constant on its support.
\item The density integrates to one.
\item For \(0<y<1\), \(p_{X\mid Y}(x\mid y)=1\) on \(0<x<y\).
\end{enumerate}
% END QUESTION Q15

\end{minipage}\par\vspace{3mm}
\par\noindent\begin{minipage}{\linewidth}
% BEGIN QUESTION Q16 / 题目与选项区
\textbf{16.} Let \(X\) take values \(-1,0,1\) with equal probabilities and let \(Y=X^2\). Select every correct claim.
\begin{enumerate}[label=(\Alph*),leftmargin=8mm,itemsep=2pt,topsep=2pt]
\item \(E[X]=0\) and \(E[Y]=2/3\).
\item \(X\) and \(Y\) are independent.
\item \(\operatorname{Corr}(X,Y)\) is undefined because \(E[X]=0\).
\item \(\operatorname{Cov}(X,Y)=2/9\).
\end{enumerate}
% END QUESTION Q16

\end{minipage}\par\vspace{3mm}
\newpage\section*{Section B. Multiple Response \hfill 24 marks}
\textit{Select all correct options. Each question has 1--3 correct options.}\par
\par\noindent\begin{minipage}{\linewidth}
% BEGIN QUESTION Q17 / 题目与选项区
\textbf{17.} The counts \((N_1,N_2,N_3)\) come from \(10\) independent categorical trials with probabilities \((0.2,0.3,0.5)\). Select every correct statement.
\begin{enumerate}[label=(\Alph*),leftmargin=8mm,itemsep=2pt,topsep=2pt]
\item Each count is Bernoulli distributed.
\item \(N_1+N_2+N_3=10\) with probability one.
\item \(N_1,N_2,N_3\) are independent because the ten trials are independent.
\item \(P(N_1=2,N_2=3,N_3=5)=\frac{10!}{2!3!5!}(0.2)^2(0.3)^3(0.5)^5\).
\end{enumerate}
% END QUESTION Q17

\end{minipage}\par\vspace{3mm}
\par\noindent\begin{minipage}{\linewidth}
% BEGIN QUESTION Q18 / 题目与选项区
\textbf{18.} For a variable with finite second moment, select every correct moment formula.
\begin{enumerate}[label=(\Alph*),leftmargin=8mm,itemsep=2pt,topsep=2pt]
\item \(\operatorname{Var}(3X-7)=9\operatorname{Var}(X)\).
\item \(E[X^2]=(E X)^2\) for every random variable.
\item \(\operatorname{Cov}(2X+1,-Y+4)=-2\operatorname{Cov}(X,Y)\).
\item \(E[(X-E X)^2]=E[X^2]-(E X)^2\).
\end{enumerate}
% END QUESTION Q18

\end{minipage}\par\vspace{3mm}
\par\noindent\begin{minipage}{\linewidth}
% BEGIN QUESTION Q19 / 题目与选项区
\textbf{19.} A nonsingular bivariate Gaussian has mean \((1,-2)^T\) and covariance \(\Sigma=\begin{pmatrix}4&0\\0&9\end{pmatrix}\). Select every correct statement.
\begin{enumerate}[label=(\Alph*),leftmargin=8mm,itemsep=2pt,topsep=2pt]
\item The two coordinates are independent.
\item The exponent contains \((x_1-1)^2/4+(x_2+2)^2/9\), multiplied by \(-1/2\).
\item The density normalization contains \(|\Sigma|^{1/2}=6\).
\item The density at its mean equals the probability that the vector equals its mean.
\end{enumerate}
% END QUESTION Q19

\end{minipage}\par\vspace{3mm}
\par\noindent\begin{minipage}{\linewidth}
% BEGIN QUESTION Q20 / 题目与选项区
\textbf{20.} A count \(K\) records successes in \(20\) independent Bernoulli trials with success probability \(0.15\). Select every correct formula.
\begin{enumerate}[label=(\Alph*),leftmargin=8mm,itemsep=2pt,topsep=2pt]
\item The fraction \(K/20\) has variance \(2.55/20\).
\item \(P(K=2)=\binom{20}{2}(0.15)^2(0.85)^{18}\).
\item \(E[K]=3\).
\item \(\operatorname{Var}(K)=3\).
\end{enumerate}
% END QUESTION Q20

\end{minipage}\par\vspace{3mm}
\newpage\section*{Section B. Multiple Response \hfill 24 marks}
\textit{Select all correct options. Each question has 1--3 correct options.}\par
\par\noindent\begin{minipage}{\linewidth}
% BEGIN QUESTION Q21 / 题目与选项区
\textbf{21.} For iid \(X_i\sim N(\mu,\sigma^2)\), \(n\ge3\), let \(S^2=\sum_i(X_i-\bar X)^2/(n-1)\). Select every correct distribution statement.
\begin{enumerate}[label=(\Alph*),leftmargin=8mm,itemsep=2pt,topsep=2pt]
\item \(\bar X\) and \(S^2\) are independent.
\item \(\sqrt n(\bar X-\mu)/S\sim t_n\).
\item \(nS^2/\sigma^2\sim\chi^2_n\).
\item \(\sqrt n(\bar X-\mu)/\sigma\sim N(0,1)\).
\end{enumerate}
% END QUESTION Q21

\end{minipage}\par\vspace{3mm}
\par\noindent\begin{minipage}{\linewidth}
% BEGIN QUESTION Q22 / 题目与选项区
\textbf{22.} Let \(T\sim t_5\) and \(W=4+2T\). Select every correct claim.
\begin{enumerate}[label=(\Alph*),leftmargin=8mm,itemsep=2pt,topsep=2pt]
\item \(T\) is exactly standard normal because its degrees of freedom exceed two.
\item \(W\) has a general t distribution with location \(4\) and scale \(2\).
\item \(\operatorname{Var}(W)=4\).
\item \(E[W]=4\).
\end{enumerate}
% END QUESTION Q22

\end{minipage}\par\vspace{3mm}
\par\noindent\begin{minipage}{\linewidth}
% BEGIN QUESTION Q23 / 题目与选项区
\textbf{23.} Two independent fair bits \(U,V\) are drawn. Let \(A=\{U=1\}\), \(B=\{V=1\}\), and \(C=\{U=V\}\). Select every correct claim.
\begin{enumerate}[label=(\Alph*),leftmargin=8mm,itemsep=2pt,topsep=2pt]
\item \(P(C\mid A\cap B)=1\).
\item Pairwise independence always implies independence of every triple intersection.
\item \(A\) and \(C\) are independent.
\item \(A,B,C\) are mutually independent.
\end{enumerate}
% END QUESTION Q23

\end{minipage}\par\vspace{3mm}
\par\noindent\begin{minipage}{\linewidth}
% BEGIN QUESTION Q24 / 题目与选项区
\textbf{24.} A binary joint distribution is \(p(0,0)=0.30\), \(p(0,1)=0.10\), \(p(1,0)=0.20\), \(p(1,1)=0.40\). Select every correct conclusion.
\begin{enumerate}[label=(\Alph*),leftmargin=8mm,itemsep=2pt,topsep=2pt]
\item \(P(X=1)=0.60\).
\item \(P(X=0\mid Y=0)=0.30\).
\item \(X\) and \(Y\) are independent.
\item \(P(Y=1\mid X=1)=2/3\).
\end{enumerate}
% END QUESTION Q24

\end{minipage}\par\vspace{3mm}
\newpage\section*{Section C. Proofs and Formula Derivations}
\textbf{25. Deriving total probability, Bayes' rule, and posterior odds}\hfill\textbf{[10 marks]}\par\medskip
\begin{enumerate}[label=(\alph*),leftmargin=8mm,itemsep=3pt]
\item Let \(A_1,\ldots,A_k\) be a finite partition with \(P(A_i)>0\). Starting from disjoint events and the product rule, derive a formula for \(P(B)\). State the additional condition needed to condition on \(B\).\hfill\textbf{[3 marks]}
\item Use the result to derive \(P(A_i\mid B)\), and prove that these posterior probabilities sum to one.\hfill\textbf{[3 marks]}
\item For two hypotheses \(H_1,H_0\) with positive priors, derive posterior odds after observing events \(E\) and \(F\). Assume the relevant likelihoods are positive. Show when the evidence contribution factors into \(P(E\mid H_i)P(F\mid H_i)\).\hfill\textbf{[4 marks]}
\end{enumerate}\par\medskip\hrule\medskip\textbf{Student response}\par
% BEGIN ANSWER Q25 / 第25题答题区
\begin{enumerate}[label=(\alph*),leftmargin=8mm,itemsep=3pt]
\item
% >>> BEGIN ANSWER Q25(a)
% 在这里写英文推导或证明；完成后可删除下方空白占位命令。
\AnswerBlank
% <<< END ANSWER Q25(a)
\item
% >>> BEGIN ANSWER Q25(b)
% 在这里写英文推导或证明；完成后可删除下方空白占位命令。
\AnswerBlank
% <<< END ANSWER Q25(b)
\item
% >>> BEGIN ANSWER Q25(c)
% 在这里写英文推导或证明；完成后可删除下方空白占位命令。
\AnswerBlank
% <<< END ANSWER Q25(c)
\end{enumerate}
% END ANSWER Q25
\newpage\section*{Section C. Proofs and Formula Derivations}
\textbf{26. Variance formulas and the effect of shared noise}\hfill\textbf{[10 marks]}\par\medskip
\begin{enumerate}[label=(\alph*),leftmargin=8mm,itemsep=3pt]
\item For finite-second-moment variables, derive \(\operatorname{Cov}(X,Y)=E[XY]-E[X]E[Y]\) and the variance of \(aX+bY+c\).\hfill\textbf{[3 marks]}
\item Let \(X_1,\ldots,X_n\) have common mean \(\mu\) and finite variance \(\sigma^2\). Derive the mean and variance of \(\bar X_n\) in terms of all pairwise covariances. Specialize to independent observations.\hfill\textbf{[4 marks]}
\item Suppose measurements have the form \(X_i=\mu+B+\epsilon_i\), where \(B,\epsilon_1,\ldots,\epsilon_n\) are independent, centered, and have variances \(\tau^2,\sigma_\epsilon^2,\ldots,\sigma_\epsilon^2\). Derive \(\operatorname{Var}(\bar X_n)\) and explain why the iid variance formula would fail.\hfill\textbf{[3 marks]}
\end{enumerate}\par\medskip\hrule\medskip\textbf{Student response}\par
% BEGIN ANSWER Q26 / 第26题答题区
\begin{enumerate}[label=(\alph*),leftmargin=8mm,itemsep=3pt]
\item
% >>> BEGIN ANSWER Q26(a)
% 在这里写英文推导或证明；完成后可删除下方空白占位命令。
\AnswerBlank
% <<< END ANSWER Q26(a)
\item
% >>> BEGIN ANSWER Q26(b)
% 在这里写英文推导或证明；完成后可删除下方空白占位命令。
\AnswerBlank
% <<< END ANSWER Q26(b)
\item
% >>> BEGIN ANSWER Q26(c)
% 在这里写英文推导或证明；完成后可删除下方空白占位命令。
\AnswerBlank
% <<< END ANSWER Q26(c)
\end{enumerate}
% END ANSWER Q26
\newpage\section*{Section C. Proofs and Formula Derivations}
\textbf{27. Constructing the binomial distribution from trials}\hfill\textbf{[10 marks]}\par\medskip
\begin{enumerate}[label=(\alph*),leftmargin=8mm,itemsep=3pt]
\item Let \(I_1,\ldots,I_n\) be independent Bernoulli variables with common success probability \(p\in(0,1)\). For \(K=\sum_i I_i\), derive \(P(K=k)\) by counting disjoint outcomes.\hfill\textbf{[3 marks]}
\item Derive the mean and variance of one Bernoulli variable directly from its two possible values. Then derive \(E[K]\) without differentiating a generating function.\hfill\textbf{[3 marks]}
\item Derive \(\operatorname{Var}(K)\) and \(\operatorname{Var}(K/n)\). Identify precisely where independence is used, and state the variance formula if the indicators are dependent but retain the same marginal success probability.\hfill\textbf{[4 marks]}
\end{enumerate}\par\medskip\hrule\medskip\textbf{Student response}\par
% BEGIN ANSWER Q27 / 第27题答题区
\begin{enumerate}[label=(\alph*),leftmargin=8mm,itemsep=3pt]
\item
% >>> BEGIN ANSWER Q27(a)
% 在这里写英文推导或证明；完成后可删除下方空白占位命令。
\AnswerBlank
% <<< END ANSWER Q27(a)
\item
% >>> BEGIN ANSWER Q27(b)
% 在这里写英文推导或证明；完成后可删除下方空白占位命令。
\AnswerBlank
% <<< END ANSWER Q27(b)
\item
% >>> BEGIN ANSWER Q27(c)
% 在这里写英文推导或证明；完成后可删除下方空白占位命令。
\AnswerBlank
% <<< END ANSWER Q27(c)
\end{enumerate}
% END ANSWER Q27
\newpage\section*{Section C. Proofs and Formula Derivations}
\textbf{28. An exact Gaussian mean and residual decomposition}\hfill\textbf{[10 marks]}\par\medskip
\begin{enumerate}[label=(\alph*),leftmargin=8mm,itemsep=3pt]
\item Let \(X_i\) be iid \(N(\mu,\sigma^2)\), \(\sigma>0\), \(n\ge2\), and set \(Z_i=(X_i-\mu)/\sigma\). Choose an orthogonal matrix \(A\) whose first row is \((1,\ldots,1)/\sqrt n\), and let \(Y=AZ\). Using the joint standard normal density, prove that the coordinates of \(Y\) are independent standard normals.\hfill\textbf{[3 marks]}
\item Express \(\sqrt n(\bar X-\mu)/\sigma\) and \(\sum_i(X_i-\bar X)^2/\sigma^2\) using these coordinates. Derive the chi-square law and prove independence of the sample mean and sample variance.\hfill\textbf{[4 marks]}
\item For \(S^2=\sum_i(X_i-\bar X)^2/(n-1)\), derive the exact distribution of \(\sqrt n(\bar X-\mu)/S\). Explain why the same exact assertion does not follow merely from iid sampling with finite variance.\hfill\textbf{[3 marks]}
\end{enumerate}\par\medskip\hrule\medskip\textbf{Student response}\par
% BEGIN ANSWER Q28 / 第28题答题区
\begin{enumerate}[label=(\alph*),leftmargin=8mm,itemsep=3pt]
\item
% >>> BEGIN ANSWER Q28(a)
% 在这里写英文推导或证明；完成后可删除下方空白占位命令。
\AnswerBlank
% <<< END ANSWER Q28(a)
\item
% >>> BEGIN ANSWER Q28(b)
% 在这里写英文推导或证明；完成后可删除下方空白占位命令。
\AnswerBlank
% <<< END ANSWER Q28(b)
\item
% >>> BEGIN ANSWER Q28(c)
% 在这里写英文推导或证明；完成后可删除下方空白占位命令。
\AnswerBlank
% <<< END ANSWER Q28(c)
\end{enumerate}
% END ANSWER Q28
\newpage\section*{Section D. Applications}
\textbf{29. Two inspections and the base rate of defective components}\hfill\textbf{[12 marks]}\par\medskip
\begin{enumerate}[label=(\alph*),leftmargin=8mm,itemsep=3pt]
\item A manufacturer draws components from a stable production stream. A component is defective with probability \(0.03\). One inspection flags a defective component with probability \(0.90\), and flags a nondefective component with probability \(0.04\). Compute the flag probability and the probability a flagged component is defective.\hfill\textbf{[4 marks]}
\item The component is inspected twice. Assume the inspection outcomes are independent conditional on whether that component is defective, with the same sensitivity and false-flag probability each time. Compute the probability of two flags and the defect probability given two flags. Explain the factorization you used.\hfill\textbf{[4 marks]}
\item A separate batch contains \(200\) independently drawn components, each inspected once, with independent inspection noise across components. Let \(K\) be its flag count. State its exact distribution, mean, variance, and the variance of its flag fraction \(K/200\).\hfill\textbf{[4 marks]}
\end{enumerate}\par\medskip\hrule\medskip\textbf{Student response}\par
% BEGIN ANSWER Q29 / 第29题答题区
\begin{enumerate}[label=(\alph*),leftmargin=8mm,itemsep=3pt]
\item
% >>> BEGIN ANSWER Q29(a)
% 在这里写英文推导或证明；完成后可删除下方空白占位命令。
\AnswerBlank
% <<< END ANSWER Q29(a)
\item
% >>> BEGIN ANSWER Q29(b)
% 在这里写英文推导或证明；完成后可删除下方空白占位命令。
\AnswerBlank
% <<< END ANSWER Q29(b)
\item
% >>> BEGIN ANSWER Q29(c)
% 在这里写英文推导或证明；完成后可删除下方空白占位命令。
\AnswerBlank
% <<< END ANSWER Q29(c)
\end{enumerate}
% END ANSWER Q29
\newpage\section*{Section D. Applications}
\textbf{30. Exact inference for a normally distributed calibration error}\hfill\textbf{[12 marks]}\par\medskip
\begin{enumerate}[label=(\alph*),leftmargin=8mm,itemsep=3pt]
\item Nine independent calibration measurements, in millimeters, are \[498,\ 502,\ 501,\ 499,\ 500,\ 503,\ 497,\ 500,\ 500.\] Assume an iid \(N(\mu,\sigma^2)\) model with both parameters unknown. Compute \(\bar x\), the residual sum of squares, the unbiased sample variance \(s^2\), and the estimated standard error of the mean.\hfill\textbf{[4 marks]}
\item Compute the t statistic for the target mean \(\mu_0=499\). Using \(t_{8,0.975}=2.306\), derive and calculate a central \(95\%\) confidence interval for \(\mu\). State the exact sampling assumption behind this interval.\hfill\textbf{[4 marks]}
\item For the reference variance \(\sigma_0^2=4\), compute the chi-square statistic. Use \(\chi^2_{8,0.025}=2.180\) and \(\chi^2_{8,0.975}=17.535\) to derive and calculate a central \(95\%\) confidence interval for \(\sigma^2\).\hfill\textbf{[4 marks]}
\end{enumerate}\par\medskip\hrule\medskip\textbf{Student response}\par
% BEGIN ANSWER Q30 / 第30题答题区
\begin{enumerate}[label=(\alph*),leftmargin=8mm,itemsep=3pt]
\item
% >>> BEGIN ANSWER Q30(a)
% 在这里写英文推导或证明；完成后可删除下方空白占位命令。
\AnswerBlank
% <<< END ANSWER Q30(a)
\item
% >>> BEGIN ANSWER Q30(b)
% 在这里写英文推导或证明；完成后可删除下方空白占位命令。
\AnswerBlank
% <<< END ANSWER Q30(b)
\item
% >>> BEGIN ANSWER Q30(c)
% 在这里写英文推导或证明；完成后可删除下方空白占位命令。
\AnswerBlank
% <<< END ANSWER Q30(c)
\end{enumerate}
% END ANSWER Q30
\end{document}
